\scanpage{69}

\mspage{2}

alors $\mathbf{\Gamma}$ opère non seulement sur \deleted{$X$} l'ens. $X$, mais aussi sur le groupoïde $\widehat\Pi_1(X)$, à cause du fait que celui-ci admet une définition purement schématique, ne faisant p intervenir la topologie de $\mathbb C$. Bien entendu, si $\lambda:\deleted{\unsure{x}}x\to y$ dans $\Pi_1$, alors pour $u\in\Gamma$, $u(\lambda)$ n'est p en général une flèche de $\Pi_1$ (de $u(x)$ vers $u(y)$) mais seulement de $\widehat\Pi_1$.

\deleted{\unsure{Notation}} Notation ${}^ux$ ($x\in X,u\in\mathbf{\Gamma}$) et \deleted{\unsure{``}}

\[
\begin{array}{ll}
{}^u\lambda:{}^ux-{}^uy&\text{si }\lambda:x\to y.\\[-.4ex]
&(\text{dans }\widehat\Pi_1(X))
\end{array}
\]

On a

\diagram{\(
(6)\qquad\left\{\begin{aligned}
\unsure{{}^u(\operatorname{id}_x)|u(\operatorname{id}_x)}=\operatorname{id}_{{}^ux}&\qquad(x\in X)\\[1.2ex]
\unsure{u(\mu\lambda)|{}^u(\mu\lambda)}={}^u\mu\,{}^u\lambda&\qquad\smash{\raisebox{-1.3em}{$\left(\begin{tikzcd}[ampersand replacement=\&,column sep=2em,row sep=1.8em]
x\arrow[r,"\lambda"]\&y\arrow[r,"\mu"]\&z\\
{}^ux\arrow[r,"{}^u\lambda"]\&{}^uy\arrow[r,"{}^u\mu"]\&{}^uz
\end{tikzcd}\right)$}}
\end{aligned}\right.
\)\vspace{2.5em}}

Notre objet est d'essayer de ``déterminer'' autant que faire se peut cette action de $\mathbf{\Gamma}$. À vrai dire, si $\bar K$ est la clôture algébrique de $K$ dans $\mathbb C$, on s'intéressera surtout aux sous-groupoïdes de $\Pi_1(X)$ et $\widehat\Pi_1(X)$ induits sur l'ens d'objets $X_{\bar K}$, qu'on désignera par $\Pi_1(X_{\bar K}),\widehat\Pi_1(X_{\bar K})$ -- le deuxième est stable par l'action de $\mathbf{\Gamma}$, et si on pose

\[
(7)\qquad\Gamma=\operatorname{Aut}_{K\text{-}\unsure{\mathrm{cat}|\mathrm{cor}}}(\bar K)=\operatorname{Gal}_{\bar K/K},
\]

de sorte que $\Gamma$ est un quotient de $\mathbf{\Gamma}$, l'action de $\mathbf{\Gamma}$ sur $\Pi_1(X_{\bar K})$ se fait à travers $\Gamma$. $\uparrow$

\margin{68}
